%=============================================================================!
% 6.1  COORDINATES THAT FIT                                                   !
%=============================================================================!
\section{Coordinates that fit the problem}
\label{sec:la-coordinates}

The bead of \cref{sec:en-track} slides without friction along a wire bent
into the curve $y = h(x)$. Its position has two coordinates, $x$ and $y$,
but the wire allows only one of them to be chosen: given $x$, the bead is
at height $h(x)$. This section describes such motions by the coordinates
they actually have, and writes their energies in those coordinates.

\subsection*{The bead and the wire}

Newton's second law for the bead needs every force on it, and one of them,
the push of the wire, is unknown. It is whatever force keeps the bead on
the wire, at right angles to the wire, and its size depends on the part
of the weight across the wire as well as on how fast the bead goes and how
sharply the wire bends. \Cref{sec:en-track}
sidestepped it with energy: the push does no work, so the speed at each
height follows from $\frac12 mv^2 + mgh = E$. That gave the speed at every
point, but not directly when the bead gets there.

The energy is simplest in the one coordinate that the wire leaves free,
$x$. The bead is at $(x, h(x))$, so by the chain rule its velocity is
\begin{equation*}
   \vb{v} = \Bigl(\dot{x},\ \frac{\dd h}{\dd x}\,\dot{x}\Bigr)
   = \dot{x}\,\bigl(1,\ h'(x)\bigr) ,
\end{equation*}
the dot marking the rate of change with time (\cref{sec:mo-acceleration})
and $h'$ the slope of the wire. The square of its length, the speed
squared, is $\dot{x}^2 + h'(x)^2\dot{x}^2 = \dot{x}^2\,(1 + h'^2)$ by
Pythagoras' theorem, so
\begin{equation*}
   K = \tfrac12 m\,\bigl(1 + h'(x)^2\bigr)\,\dot{x}^2 ,
   \qquad U = mg\,h(x) .
\end{equation*}
Both energies are now written with $x$ and $\dot{x}$ alone, and the push
of the wire appears in neither.

\subsection*{The pendulum}

A pendulum's bob on a light rod of length $L$ (\cref{sec:os-small}) is
fixed by the angle $\theta$ of the rod from the vertical. With the pivot
at the origin and $y$ up, the bob is at $(L\sin\theta, -L\cos\theta)$.
Differentiating each component with the chain rule, $\dd\sin\theta/\dd t =
\cos\theta\,\dot\theta$ and $\dd\cos\theta/\dd t = -\sin\theta\,\dot\theta$
(\cref{sec:ro-plane}),
\begin{equation*}
   \vb{v} = \bigl(L\cos\theta\,\dot\theta,\ L\sin\theta\,\dot\theta\bigr)
   = L\dot\theta\,(\cos\theta, \sin\theta) ,
\end{equation*}
of length $L\lvert\dot\theta\rvert$ by Pythagoras' theorem, since
$\cos^2\theta + \sin^2\theta = 1$. The bob is $L\cos\theta$ below the
pivot, so $L - L\cos\theta$ above the bottom of the swing
(\cref{sec:os-small}), and
\begin{equation*}
   K = \tfrac12 mL^2\dot\theta^2 ,
   \qquad U = mgL\,(1 - \cos\theta) .
\end{equation*}

\subsection*{Generalised coordinates}

The fewest numbers that fix the position of every part of a system, once
its constraints are obeyed, are called its generalised coordinates; each
is free to change whatever the others do, and we write them $q_1, q_2,
\dots, q_n$. Their rates of change, $\dot{q}_1,
\dots, \dot{q}_n$, are the generalised velocities. The bead has one, $x$;
the pendulum one, $\theta$; the puck of \cref{sec:ro-plane} on its table
two, $r$ and $\theta$; and $N$ free atoms have $3N$, their Cartesian
coordinates. The number $n$ is the number of degrees of freedom of
\cref{sec:ro-freedom}. A constraint that fixes one coordinate in terms of
the others, as the wire fixes $y = h(x)$ or the rod fixes the distance
of the bob from the pivot, is called holonomic, and it removes one degree
of freedom. The wire and the rod do no work, since they push at right
angles to every motion they allow; all the constraints in this book are
holonomic and do no work.

\begin{keyidea}
Generalised coordinates are the fewest numbers that fix the position of
a system once its constraints are obeyed. In them the kinetic and potential
energies can be written without the forces of the constraints, which do
no work.
\end{keyidea}

\begin{selfcheck}
How many generalised coordinates does a double pendulum have, one bob hanging
by a rod from a second bob that hangs by a rod from a fixed pivot, moving
in a plane, and how many a rigid diatomic molecule free in space?
\end{selfcheck}
